\documentclass{article}

\usepackage{amsmath, amssymb, amsthm}
\usepackage[a4paper, margin=1in]{geometry}

\title{Euclid's Algorithm}
\author{sid}
\date{\today}

\newtheorem{theorem}{Theorem}
\newtheorem{lemma}{Lemma}
\newtheorem{example}{Example}

\newcommand{\Z}{\mathbb{Z}}
\newcommand{\bld}[1]{\textbf{#1}}
\newcommand{\itl}[1]{\textit{#1}}
\newcommand{\N}{\mathbb{N}}
\newcommand{\Q}{\mathbb{Q}}
\newcommand{\R}{\mathbb{R}}
\newcommand{\eps}{\varepsilon}
\newcommand{\poly}{\mathrm{poly}}
\newcommand{\polylog}{\mathrm{polylog}}
\newcommand{\bigO}{\mathcal{O}}

\begin{document}

\maketitle

\section{Goal}
    Given two integers $a$ and $b$, compute
    $$ gcd(a,b) $$

\section{Core Idea}
    $$ gcd(a,b)=gcd(b,a\bmod{b}) $$

\section{Za Algorithm}
    Given integers $a > b$:

    \begin{enumerate}
        \item Compute $r = a \bmod b$
        \item Replace $a \leftarrow b$, $b \leftarrow r$
        \item Repeat until $b = 0$
        \item Answer is $a$
    \end{enumerate}

    \begin{example}
        \[
        \begin{aligned}
        \gcd(48, 18)
        &= \gcd(18, 48 \bmod 18) \\
        &= \gcd(18, 12) \\
        &= \gcd(12, 6) \\
        &= \gcd(6, 0) \\
        &= 6
        \end{aligned}
        \]
    \end{example}

\subsection*{Why Does This Work?}

\begin{proof}
Let $r = a \bmod b$. By the division algorithm,
\[
a = qb + r
\]
for some integer $q$.

We show that the set of common divisors of $(a,b)$ is equal to that of $(b,r)$.

\medskip
\noindent\textbf{Forward direction.}
Suppose $d \mid a$ and $d \mid b$. Then:
\[
d \mid (a - qb).
\]
Since $a - qb = r$, we get:
\[
d \mid r.
\]
Thus, $d$ divides both $b$ and $r$.

\medskip
\noindent\textbf{Reverse direction.}
Suppose $d \mid b$ and $d \mid r$. Then:
\[
d \mid (qb + r).
\]
Since $qb + r = a$, we get:
\[
d \mid a.
\]
Thus, $d$ divides both $a$ and $b$.

\medskip
Hence, the set of common divisors of $(a,b)$ and $(b,r)$ are identical, and therefore:
\[
\gcd(a,b) = \gcd(b,r).
\]
\end{proof}
\end{document}